\documentclass[11pt,fleqn]{book} % Default font size and left-justified equations

\usepackage{graphicx} % Required for including pictures
\graphicspath{{Pictures/}} % Specifies the directory where pictures are stored

\usepackage{lipsum} % Inserts dummy text

\usepackage{tikz} % Required for drawing custom shapes

\usepackage[brazil]{babel} % English language/hyphenation

\usepackage{enumitem} % Customize lists
\setlist{nolistsep} % Reduce spacing between bullet points and numbered lists

\usepackage{booktabs} % Required for nicer horizontal rules in tables

\usepackage{xcolor} % Required for specifying colors by name
\definecolor{ocre}{RGB}{116,23,255} % Define the orange color used for highlighting throughout the book

%----------------------------------------------------------------------------------------
%	MARGINS
%----------------------------------------------------------------------------------------

\usepackage{geometry} % Required for adjusting page dimensions and margins

\geometry{
	paper=a4paper, % Paper size, change to letterpaper for US letter size
	top=2cm, % Top margin
	bottom=2cm, % Bottom margin
	left=3cm, % Left margin
	right=3cm, % Right margin
	headheight=14pt, % Header height
	footskip=1.4cm, % Space from the bottom margin to the baseline of the footer
	headsep=10pt, % Space from the top margin to the baseline of the header
	%showframe, % Uncomment to show how the type block is set on the page
}

%----------------------------------------------------------------------------------------
%	FONTS
%----------------------------------------------------------------------------------------

\usepackage{avant} % Use the Avantgarde font for headings
%\usepackage{times} % Use the Times font for headings
\usepackage{mathptmx} % Use the Adobe Times Roman as the default text font together with math symbols from the Sym­bol, Chancery and Com­puter Modern fonts

\usepackage{microtype} % Slightly tweak font spacing for aesthetics
\usepackage[utf8]{inputenc} % Required for including letters with accents
\usepackage[T1]{fontenc} % Use 8-bit encoding that has 256 glyphs

%----------------------------------------------------------------------------------------
%	BIBLIOGRAPHY AND INDEX
%----------------------------------------------------------------------------------------

\usepackage[style=numeric,citestyle=numeric,sorting=nyt,sortcites=true,autopunct=true,babel=hyphen,hyperref=true,abbreviate=false,backref=true,backend=biber]{biblatex}
\addbibresource{bibliography.bib} % BibTeX bibliography file
\defbibheading{bibempty}{}

\usepackage{calc} % For simpler calculation - used for spacing the index letter headings correctly
\usepackage{makeidx} % Required to make an index
\makeindex % Tells LaTeX to create the files required for indexing

%----------------------------------------------------------------------------------------
%	MAIN TABLE OF CONTENTS
%----------------------------------------------------------------------------------------

\usepackage{titletoc} % Required for manipulating the table of contents

\contentsmargin{0cm} % Removes the default margin

% Part text styling (this is mostly taken care of in the PART HEADINGS section of this file)
\titlecontents{part}
	[0cm] % Left indentation
	{\addvspace{20pt}\bfseries} % Spacing and font options for parts
	{}
	{}
	{}

% Chapter text styling
\titlecontents{chapter}
	[1.25cm] % Left indentation
	{\addvspace{12pt}\large\sffamily\bfseries} % Spacing and font options for chapters
	{\color{ocre!60}\contentslabel[\Large\thecontentslabel]{1.25cm}\color{ocre}} % Formatting of numbered sections of this type
	{\color{ocre}} % Formatting of numberless sections of this type
	{\color{ocre!60}\normalsize\;\titlerule*[.5pc]{.}\;\thecontentspage} % Formatting of the filler to the right of the heading and the page number

% Section text styling
\titlecontents{section}
	[1.25cm] % Left indentation
	{\addvspace{3pt}\sffamily\bfseries} % Spacing and font options for sections
	{\contentslabel[\thecontentslabel]{1.25cm}} % Formatting of numbered sections of this type
	{} % Formatting of numberless sections of this type
	{\hfill\color{black}\thecontentspage} % Formatting of the filler to the right of the heading and the page number

% Subsection text styling
\titlecontents{subsection}
	[1.25cm] % Left indentation
	{\addvspace{1pt}\sffamily\small} % Spacing and font options for subsections
	{\contentslabel[\thecontentslabel]{1.25cm}} % Formatting of numbered sections of this type
	{} % Formatting of numberless sections of this type
	{\ \titlerule*[.5pc]{.}\;\thecontentspage} % Formatting of the filler to the right of the heading and the page number

% Figure text styling
\titlecontents{figure}
	[1.25cm] % Left indentation
	{\addvspace{1pt}\sffamily\small} % Spacing and font options for figures
	{\thecontentslabel\hspace*{1em}} % Formatting of numbered sections of this type
	{} % Formatting of numberless sections of this type
	{\ \titlerule*[.5pc]{.}\;\thecontentspage} % Formatting of the filler to the right of the heading and the page number

% Table text styling
\titlecontents{table}
	[1.25cm] % Left indentation
	{\addvspace{1pt}\sffamily\small} % Spacing and font options for tables
	{\thecontentslabel\hspace*{1em}} % Formatting of numbered sections of this type
	{} % Formatting of numberless sections of this type
	{\ \titlerule*[.5pc]{.}\;\thecontentspage} % Formatting of the filler to the right of the heading and the page number

%----------------------------------------------------------------------------------------
%	MINI TABLE OF CONTENTS IN PART HEADS
%----------------------------------------------------------------------------------------

% Chapter text styling
\titlecontents{lchapter}
	[0em] % Left indentation
	{\addvspace{15pt}\large\sffamily\bfseries} % Spacing and font options for chapters
	{\color{ocre}\contentslabel[\Large\thecontentslabel]{1.25cm}\color{ocre}} % Chapter number
	{}  
	{\color{ocre}\normalsize\sffamily\bfseries\;\titlerule*[.5pc]{.}\;\thecontentspage} % Page number

% Section text styling
\titlecontents{lsection}
	[0em] % Left indentation
	{\sffamily\small} % Spacing and font options for sections
	{\contentslabel[\thecontentslabel]{1.25cm}} % Section number
	{}
	{}

% Subsection text styling (note these aren't shown by default, display them by searchings this file for tocdepth and reading the commented text)
\titlecontents{lsubsection}
	[.5em] % Left indentation
	{\sffamily\footnotesize} % Spacing and font options for subsections
	{\contentslabel[\thecontentslabel]{1.25cm}}
	{}
	{}

%----------------------------------------------------------------------------------------
%	HEADERS AND FOOTERS
%----------------------------------------------------------------------------------------

\usepackage{fancyhdr} % Required for header and footer configuration

\pagestyle{fancy} % Enable the custom headers and footers

\renewcommand{\chaptermark}[1]{\markboth{\sffamily\normalsize\bfseries\chaptername\ \thechapter.\ #1}{}} % Styling for the current chapter in the header
\renewcommand{\sectionmark}[1]{\markright{\sffamily\normalsize\thesection\hspace{5pt}#1}{}} % Styling for the current section in the header

\fancyhf{} % Clear default headers and footers
\fancyhead[LE,RO]{\sffamily\normalsize\thepage} % Styling for the page number in the header
\fancyhead[LO]{\rightmark} % Print the nearest section name on the left side of odd pages
\fancyhead[RE]{\leftmark} % Print the current chapter name on the right side of even pages
%\fancyfoot[C]{\thepage} % Uncomment to include a footer

\renewcommand{\headrulewidth}{0.5pt} % Thickness of the rule under the header

\fancypagestyle{plain}{% Style for when a plain pagestyle is specified
	\fancyhead{}\renewcommand{\headrulewidth}{0pt}%
}

% Removes the header from odd empty pages at the end of chapters
\makeatletter
\renewcommand{\cleardoublepage}{
\clearpage\ifodd\c@page\else
\hbox{}
\vspace*{\fill}
\thispagestyle{empty}
\newpage
\fi}

%----------------------------------------------------------------------------------------
%	THEOREM STYLES
%----------------------------------------------------------------------------------------

\usepackage{amsmath,amsfonts,amssymb,amsthm} % For math equations, theorems, symbols, etc

\newcommand{\intoo}[2]{\mathopen{]}#1\,;#2\mathclose{[}}
\newcommand{\ud}{\mathop{\mathrm{{}d}}\mathopen{}}
\newcommand{\intff}[2]{\mathopen{[}#1\,;#2\mathclose{]}}
\renewcommand{\qedsymbol}{$\blacksquare$}
\newtheorem{notation}{Notation}[chapter]

% Boxed/framed environments
\newtheoremstyle{ocrenumbox}% Theorem style name
{0pt}% Space above
{0pt}% Space below
{\normalfont}% Body font
{}% Indent amount
{\small\bf\sffamily\color{ocre}}% Theorem head font
{\;}% Punctuation after theorem head
{0.25em}% Space after theorem head
{\small\sffamily\color{ocre}\thmname{#1}\nobreakspace\thmnumber{\@ifnotempty{#1}{}\@upn{#2}}% Theorem text (e.g. Theorem 2.1)
\thmnote{\nobreakspace\the\thm@notefont\sffamily\bfseries\color{black}---\nobreakspace#3.}} % Optional theorem note

\newtheoremstyle{blacknumex}% Theorem style name
{5pt}% Space above
{5pt}% Space below
{\normalfont}% Body font
{} % Indent amount
{\small\bf\sffamily}% Theorem head font
{\;}% Punctuation after theorem head
{0.25em}% Space after theorem head
{\small\sffamily{\tiny\ensuremath{\blacksquare}}\nobreakspace\thmname{#1}\nobreakspace\thmnumber{\@ifnotempty{#1}{}\@upn{#2}}% Theorem text (e.g. Theorem 2.1)
\thmnote{\nobreakspace\the\thm@notefont\sffamily\bfseries---\nobreakspace#3.}}% Optional theorem note

\newtheoremstyle{blacknumbox} % Theorem style name
{0pt}% Space above
{0pt}% Space below
{\normalfont}% Body font
{}% Indent amount
{\small\bf\sffamily}% Theorem head font
{\;}% Punctuation after theorem head
{0.25em}% Space after theorem head
{\small\sffamily\thmname{#1}\nobreakspace\thmnumber{\@ifnotempty{#1}{}\@upn{#2}}% Theorem text (e.g. Theorem 2.1)
\thmnote{\nobreakspace\the\thm@notefont\sffamily\bfseries---\nobreakspace#3.}}% Optional theorem note

% Non-boxed/non-framed environments
\newtheoremstyle{ocrenum}% Theorem style name
{5pt}% Space above
{5pt}% Space below
{\normalfont}% Body font
{}% Indent amount
{\small\bf\sffamily\color{ocre}}% Theorem head font
{\;}% Punctuation after theorem head
{0.25em}% Space after theorem head
{\small\sffamily\color{ocre}\thmname{#1}\nobreakspace\thmnumber{\@ifnotempty{#1}{}\@upn{#2}}% Theorem text (e.g. Theorem 2.1)
\thmnote{\nobreakspace\the\thm@notefont\sffamily\bfseries\color{black}---\nobreakspace#3.}} % Optional theorem note
\makeatother

% Defines the theorem text style for each type of theorem to one of the three styles above
\newcounter{dummy} 
\numberwithin{dummy}{section}
\theoremstyle{ocrenumbox}
\newtheorem{theoremeT}[dummy]{Theorem}
\newtheorem{problem}{Problema}[chapter]
\newtheorem{exerciseT}{Exercício}[chapter]
\theoremstyle{blacknumex}
\newtheorem{exampleT}{Exemplo}[chapter]
\theoremstyle{blacknumbox}
\newtheorem{vocabulary}{Vocabulary}[chapter]
\newtheorem{definitionT}{Definição}[section]
\newtheorem{corollaryT}[dummy]{Corollary}
\theoremstyle{ocrenum}
\newtheorem{proposition}[dummy]{Proposition}

%----------------------------------------------------------------------------------------
%	DEFINITION OF COLORED BOXES
%----------------------------------------------------------------------------------------

\RequirePackage[framemethod=default]{mdframed} % Required for creating the theorem, definition, exercise and corollary boxes

% Theorem box
\newmdenv[skipabove=7pt,
skipbelow=7pt,
backgroundcolor=black!5,
linecolor=ocre,
innerleftmargin=5pt,
innerrightmargin=5pt,
innertopmargin=5pt,
leftmargin=0cm,
rightmargin=0cm,
innerbottommargin=5pt]{tBox}

% Exercise box	  
\newmdenv[skipabove=7pt,
skipbelow=7pt,
rightline=false,
leftline=true,
topline=false,
bottomline=false,
backgroundcolor=ocre!10,
linecolor=ocre,
innerleftmargin=5pt,
innerrightmargin=5pt,
innertopmargin=5pt,
innerbottommargin=5pt,
leftmargin=0cm,
rightmargin=0cm,
linewidth=4pt]{eBox}	

% Definition box
\newmdenv[skipabove=7pt,
skipbelow=7pt,
rightline=false,
leftline=true,
topline=false,
bottomline=false,
linecolor=ocre,
innerleftmargin=5pt,
innerrightmargin=5pt,
innertopmargin=0pt,
leftmargin=0cm,
rightmargin=0cm,
linewidth=4pt,
innerbottommargin=0pt]{dBox}	

% Corollary box
\newmdenv[skipabove=7pt,
skipbelow=7pt,
rightline=false,
leftline=true,
topline=false,
bottomline=false,
linecolor=gray,
backgroundcolor=black!5,
innerleftmargin=5pt,
innerrightmargin=5pt,
innertopmargin=5pt,
leftmargin=0cm,
rightmargin=0cm,
linewidth=4pt,
innerbottommargin=5pt]{cBox}

% Creates an environment for each type of theorem and assigns it a theorem text style from the "Theorem Styles" section above and a colored box from above
\newenvironment{theorem}{\begin{tBox}\begin{theoremeT}}{\end{theoremeT}\end{tBox}}
\newenvironment{exercise}{\begin{eBox}\begin{exerciseT}}{\hfill{\color{ocre}\tiny\ensuremath{\blacksquare}}\end{exerciseT}\end{eBox}}				  
\newenvironment{definition}{\begin{dBox}\begin{definitionT}}{\end{definitionT}\end{dBox}}	
\newenvironment{example}{\begin{exampleT}}{\hfill{\tiny\ensuremath{\blacksquare}}\end{exampleT}}		
\newenvironment{corollary}{\begin{cBox}\begin{corollaryT}}{\end{corollaryT}\end{cBox}}	

%----------------------------------------------------------------------------------------
%	REMARK ENVIRONMENT
%----------------------------------------------------------------------------------------

\newenvironment{remark}{\par\vspace{10pt}\small % Vertical white space above the remark and smaller font size
\begin{list}{}{
\leftmargin=35pt % Indentation on the left
\rightmargin=25pt}\item\ignorespaces % Indentation on the right
\makebox[-2.5pt]{\begin{tikzpicture}[overlay]
\node[draw=ocre!60,line width=1pt,circle,fill=ocre!25,font=\sffamily\bfseries,inner sep=2pt,outer sep=0pt] at (-15pt,0pt){\textcolor{ocre}{R}};\end{tikzpicture}} % Orange R in a circle
\advance\baselineskip -1pt}{\end{list}\vskip5pt} % Tighter line spacing and white space after remark

%----------------------------------------------------------------------------------------
%	SECTION NUMBERING IN THE MARGIN
%----------------------------------------------------------------------------------------

\makeatletter
\renewcommand{\@seccntformat}[1]{\llap{\textcolor{ocre}{\csname the#1\endcsname}\hspace{1em}}}                    
\renewcommand{\section}{\@startsection{section}{1}{\z@}
{-4ex \@plus -1ex \@minus -.4ex}
{1ex \@plus.2ex }
{\normalfont\large\sffamily\bfseries}}
\renewcommand{\subsection}{\@startsection {subsection}{2}{\z@}
{-3ex \@plus -0.1ex \@minus -.4ex}
{0.5ex \@plus.2ex }
{\normalfont\sffamily\bfseries}}
\renewcommand{\subsubsection}{\@startsection {subsubsection}{3}{\z@}
{-2ex \@plus -0.1ex \@minus -.2ex}
{.2ex \@plus.2ex }
{\normalfont\small\sffamily\bfseries}}                        
\renewcommand\paragraph{\@startsection{paragraph}{4}{\z@}
{-2ex \@plus-.2ex \@minus .2ex}
{.1ex}
{\normalfont\small\sffamily\bfseries}}

%----------------------------------------------------------------------------------------
%	PART HEADINGS
%----------------------------------------------------------------------------------------

% Numbered part in the table of contents
\newcommand{\@mypartnumtocformat}[2]{%
	\setlength\fboxsep{0pt}%
	\noindent\colorbox{ocre!20}{\strut\parbox[c][.7cm]{\ecart}{\color{ocre!70}\Large\sffamily\bfseries\centering#1}}\hskip\esp\colorbox{ocre!40}{\strut\parbox[c][.7cm]{\linewidth-\ecart-\esp}{\Large\sffamily\centering#2}}%
}

% Unnumbered part in the table of contents
\newcommand{\@myparttocformat}[1]{%
	\setlength\fboxsep{0pt}%
	\noindent\colorbox{ocre!40}{\strut\parbox[c][.7cm]{\linewidth}{\Large\sffamily\centering#1}}%
}

\newlength\esp
\setlength\esp{4pt}
\newlength\ecart
\setlength\ecart{1.2cm-\esp}
\newcommand{\thepartimage}{}%
\newcommand{\partimage}[1]{\renewcommand{\thepartimage}{#1}}%
\def\@part[#1]#2{%
\ifnum \c@secnumdepth >-2\relax%
\refstepcounter{part}%
\addcontentsline{toc}{part}{\texorpdfstring{\protect\@mypartnumtocformat{\thepart}{#1}}{\partname~\thepart\ ---\ #1}}
\else%
\addcontentsline{toc}{part}{\texorpdfstring{\protect\@myparttocformat{#1}}{#1}}%
\fi%
\startcontents%
\markboth{}{}%
{\thispagestyle{empty}%
\begin{tikzpicture}[remember picture,overlay]%
\node at (current page.north west){\begin{tikzpicture}[remember picture,overlay]%	
\fill[ocre!20](0cm,0cm) rectangle (\paperwidth,-\paperheight);
\node[anchor=north] at (4cm,-3.25cm){\color{ocre!40}\fontsize{220}{100}\sffamily\bfseries\thepart}; 
\node[anchor=south east] at (\paperwidth-1cm,-\paperheight+1cm){\parbox[t][][t]{8.5cm}{
\printcontents{l}{0}{\setcounter{tocdepth}{1}}% The depth to which the Part mini table of contents displays headings; 0 for chapters only, 1 for chapters and sections and 2 for chapters, sections and subsections
}};
\node[anchor=north east] at (\paperwidth-1.5cm,-3.25cm){\parbox[t][][t]{15cm}{\strut\raggedleft\color{white}\fontsize{30}{30}\sffamily\bfseries#2}};
\end{tikzpicture}};
\end{tikzpicture}}%
\@endpart}
\def\@spart#1{%
\startcontents%
\phantomsection
{\thispagestyle{empty}%
\begin{tikzpicture}[remember picture,overlay]%
\node at (current page.north west){\begin{tikzpicture}[remember picture,overlay]%	
\fill[ocre!20](0cm,0cm) rectangle (\paperwidth,-\paperheight);
\node[anchor=north east] at (\paperwidth-1.5cm,-3.25cm){\parbox[t][][t]{15cm}{\strut\raggedleft\color{white}\fontsize{30}{30}\sffamily\bfseries#1}};
\end{tikzpicture}};
\end{tikzpicture}}
\addcontentsline{toc}{part}{\texorpdfstring{%
\setlength\fboxsep{0pt}%
\noindent\protect\colorbox{ocre!40}{\strut\protect\parbox[c][.7cm]{\linewidth}{\Large\sffamily\protect\centering #1\quad\mbox{}}}}{#1}}%
\@endpart}
\def\@endpart{\vfil\newpage
\if@twoside
\if@openright
\null
\thispagestyle{empty}%
\newpage
\fi
\fi
\if@tempswa
\twocolumn
\fi}

%----------------------------------------------------------------------------------------
%	CHAPTER HEADINGS
%----------------------------------------------------------------------------------------

% A switch to conditionally include a picture, implemented by Christian Hupfer
\newif\ifusechapterimage
\usechapterimagetrue
\newcommand{\thechapterimage}{}%
\newcommand{\chapterimage}[1]{\ifusechapterimage\renewcommand{\thechapterimage}{#1}\fi}%
\newcommand{\autodot}{.}
\def\@makechapterhead#1{%
{\parindent \z@ \raggedright \normalfont
\ifnum \c@secnumdepth >\m@ne
\if@mainmatter
\begin{tikzpicture}[remember picture,overlay]
\node at (current page.north west)
{\begin{tikzpicture}[remember picture,overlay]
\node[anchor=north west,inner sep=0pt] at (0,0) {\ifusechapterimage\includegraphics[width=\paperwidth]{\thechapterimage}\fi};
\draw[anchor=west] (\Gm@lmargin,-9cm) node [line width=2pt,rounded corners=15pt,draw=ocre,fill=white,fill opacity=0.5,inner sep=15pt]{\strut\makebox[22cm]{}};
\draw[anchor=west] (\Gm@lmargin+.3cm,-9cm) node {\huge\sffamily\bfseries\color{black}\thechapter\autodot~#1\strut};
\end{tikzpicture}};
\end{tikzpicture}
\else
\begin{tikzpicture}[remember picture,overlay]
\node at (current page.north west)
{\begin{tikzpicture}[remember picture,overlay]
\node[anchor=north west,inner sep=0pt] at (0,0) {\ifusechapterimage\includegraphics[width=\paperwidth]{\thechapterimage}\fi};
\draw[anchor=west] (\Gm@lmargin,-9cm) node [line width=2pt,rounded corners=15pt,draw=ocre,fill=white,fill opacity=0.5,inner sep=15pt]{\strut\makebox[22cm]{}};
\draw[anchor=west] (\Gm@lmargin+.3cm,-9cm) node {\huge\sffamily\bfseries\color{black}#1\strut};
\end{tikzpicture}};
\end{tikzpicture}
\fi\fi\par\vspace*{270\p@}}}

%-------------------------------------------

\def\@makeschapterhead#1{%
\begin{tikzpicture}[remember picture,overlay]
\node at (current page.north west)
{\begin{tikzpicture}[remember picture,overlay]
\node[anchor=north west,inner sep=0pt] at (0,0) {\ifusechapterimage\includegraphics[width=\paperwidth]{\thechapterimage}\fi};
\draw[anchor=west] (\Gm@lmargin,-9cm) node [line width=2pt,rounded corners=15pt,draw=ocre,fill=white,fill opacity=0.5,inner sep=15pt]{\strut\makebox[22cm]{}};
\draw[anchor=west] (\Gm@lmargin+.3cm,-9cm) node {\huge\sffamily\bfseries\color{black}#1\strut};
\end{tikzpicture}};
\end{tikzpicture}
\par\vspace*{270\p@}}
\makeatother

%----------------------------------------------------------------------------------------
%	LINKS
%----------------------------------------------------------------------------------------

\usepackage{hyperref}
\hypersetup{hidelinks,backref=true,pagebackref=true,hyperindex=true,colorlinks=false,breaklinks=true,urlcolor=ocre,bookmarks=true,bookmarksopen=false}

\usepackage{bookmark}
\bookmarksetup{
open,
numbered,
addtohook={%
\ifnum\bookmarkget{level}=0 % chapter
\bookmarksetup{bold}%
\fi
\ifnum\bookmarkget{level}=-1 % part
\bookmarksetup{color=ocre,bold}%
\fi
}
}
 % Insert the commands.tex file which contains the majority of the structure behind the template

%\hypersetup{pdftitle={Title},pdfauthor={Author}} % Uncomment and fill out to include PDF metadata for the author and title of the book

%----------------------------------------------------------------------------------------

\begin{document}

%----------------------------------------------------------------------------------------
%	TITLE PAGE
%----------------------------------------------------------------------------------------

\begingroup
\thispagestyle{empty} % Suppress headers and footers on the title page
\begin{tikzpicture}[remember picture,overlay]
\node[inner sep=0pt] (background) at (current page.center) {\includegraphics[width=\paperwidth]{background.pdf}};
\draw (current page.center) node [fill=ocre!30!white,fill opacity=0.6,text opacity=1,inner sep=1cm]{\Huge\centering\bfseries\sffamily\parbox[c][][t]{\paperwidth}{\centering Matemática no ENEM\\[15pt] % Book title
{\Large Monitoria AnnWay}\\[20pt] % Subtitle
{\huge Lucas, Marcelo e Matheus}}}; % Author name
\end{tikzpicture}
\vfill
\endgroup

%----------------------------------------------------------------------------------------
%	COPYRIGHT PAGE
%----------------------------------------------------------------------------------------

\newpage
~\vfill
\thispagestyle{empty}

\noindent Copyright \copyright\ 2020 Monitoria AnnWay\\ % Copyright notice

\noindent \textsc{Published by Publisher}\\ % Publisher

\noindent \textsc{http://annway.github.io/}\\ % URL

\noindent Licensed under the Creative Commons Attribution-NonCommercial 3.0 Unported License (the ``License''). You may not use this file except in compliance with the License. You may obtain a copy of the License at \url{http://creativecommons.org/licenses/by-nc/3.0}. Unless required by applicable law or agreed to in writing, software distributed under the License is distributed on an \textsc{``as is'' basis, without warranties or conditions of any kind}, either express or implied. See the License for the specific language governing permissions and limitations under the License.\\ % License information, replace this with your own license (if any)

\noindent \textit{Primeira versão, junho de 2020} % Printing/edition date

%----------------------------------------------------------------------------------------
%	TABLE OF CONTENTS
%----------------------------------------------------------------------------------------

%\usechapterimagefalse % If you don't want to include a chapter image, use this to toggle images off - it can be enabled later with \usechapterimagetrue

\chapterimage{chapter_head_1.pdf} % Table of contents heading image

\pagestyle{empty} % Disable headers and footers for the following pages

\tableofcontents % Print the table of contents itself

\cleardoublepage % Forces the first chapter to start on an odd page so it's on the right side of the book

\pagestyle{fancy} % Enable headers and footers again


\chapter{Sobre o projeto}
Sabemos que os assuntos abordados no ENEM são vastos, afinal são muitas as disciplinas que os estudantes têm contato durante a vida escolar. Mesmo uma única disciplina como história, por exemplo, cobra uma gama enorme de informação e conhecimento por parte dos estudantes. Portanto, nada melhor que organizar e priorizar seu tempo de estudo.

Por isso, Elaboramos este documento com o intuito de lhe ajudar. Elencamos os 10 assuntos que possuem maior número de questões no ENEM na área do conhecimento da Matemática e suas Aplicações. Além de você saber quais são os assuntos mais exigidos, também saberá o que é preciso pra estudá-los. A lista segue do assunto mais cobrado ao mesmo cobrado.


\part{Matemática Básica}


\chapterimage{chapter_head_2.pdf}
\chapter{Aritmética}
    \section{Operações básicas }
    \subsection{Conjuntos dos numero inteiros, naturais e racionais}


\chapter{Sistemas de Unidades de Medidas}
(LEMBRAR DE COLOCAR AS IMAGENS)
\\O Sistema Métrico é um sistema internacional de medições que determina as unidades a serem utilizadas por cada uma das medidas, além de suas transformações. A unidade fundamental desse sistema é o metro (m), porém, cada medida tem sua unidade padrão, de acordo com o Sistema Internacional de Unidades. 

\section{Medidas de comprimento}
As medidas de comprimento são os mecanismos de medição mais utilizados. A unidade fundamental das medidas de comprimento é o metro (m). 
\subsection{Múltiplos de metro}
\begin{itemize}
	\item Quilômetro(km) $= 1000m = 10^3m$
	\item Hectômetro (hm)  = 100m = $10^2$m
	\item Decâmetro (dam)= 10m 
	
\end{itemize}

\subsection{Submúltiplos de metro}
\begin{itemize}
	\item Decimetro(dm) = 0,1m = $10^{-1}m$
	\item centímetro (cm) = 0,01m = $10^{-2}m$
	\item Milímetro (mm)= 0,001m = $10^{-3}m$
\end{itemize}
\section{Medidas de Superfície}
As medidas de superfície são as utilizadas na medição de
áreas. Sua unidade fundamental é o metro quadrado ($m^2$).
OBS: COLOCAR UM QUADRADO AQUI
\subsection{Múltiplos de metro quadrado($m^2$)}
\begin{itemize}
	\item Quilômetro quadrado $(km^2) = 1.000.000 m^2 = 10^6 m^2$
	\item Hectômetro quadrado $(hm^2) = 10.000 m^2 = 10^4 m^2$
	\item Decâmetro quadrado $(dam^2) = 100 m^2 = 10^2 m^2$
	
\end{itemize}

\subsection{Submúltiplos de metro quadrado}
\begin{itemize}
	\item Decímetro quadrado $(dm^2) = 0,01m^2 = 10^{-2}m^2$
	\item centímetro quadrado $(cm^2) = 0,0001m^2 = 10^{-4}m^2$
	\item Milímetro quadrado $(mm^2)= 0,000001m^2 = 10^{-6}m^2$
\end{itemize}
\section{Medidas Agrárias}
Dentre as medidas de superfícies, existem as medidas
agrárias, que são mais utilizadas para medir áreas rurais. Sua
unidade fundamental é o are (a).

\begin{itemize}
	\item Centiare $(ca) = 1m^2=10^0m^2$
	\item Are $(a)= 100m^2 = 10^2m^2$
	\item Hectare $(ha) = 10000m^2 = 10^4m^2$
\end{itemize}
\section{Medidas de Volume} 
As medidas de volume são utilizadas g
para medir o espaço ocupado por um
corpo tridimensional ou a sua
capacidade de armazenar alguma
substância. Em química, as medidas de
volume geralmente aparecem quando
medimos quantidades de líquidos. A
unidade métrica tradicional de volume
usada nesse caso é o litro (L). Em termos
do SI, um litro é definido exatamente
como 1 decímetro cúbico. 
\subsection{Múltiplos do Metro Cubico:}
\begin{itemize}
    \item Quilômetro cúbico $(Km^3) = 1.000.000.000 m^3 = 10^9 m^3$
    \item Hectômetro cúbico $(hm^3) = 1.000.000 m^3 = 10^6 m^3$
    \item Decâmetro cúbico $(dam^3) = 1.000 m^3 = 10^3 m^3$
\end{itemize}
\subsection{Submúltiplos de metros cúbicos $(m^3)$}
\begin{itemize}
    \item Decímetro cúbico $(dm^3) = 0,001 m^3 = 10^{-3} m^3$
    \item Centímetro cúbico $(cm^3) = 0,000001 m^3 = 10^{-6} m^3$
    \item Milímetro cúbico $(mm^3) = 0,000000001 m^3 = 10^{-9} m^3$
\end{itemize}
\section{Medidas de Capacidade}
As medidas de capacidade são utilizadas para representarem
as grandezas que definem o volume contido em um certo
reservatório. A mais utilizada é o litro (l).
	\\\textbf{Importante:}
\begin{itemize}
    \item $1l = 1dm^3$
    \item $1 ml = 1 cm^3$
    \item $10^3 l = 1 m^3$
\end{itemize}
\subsection{Múltiplos de litro:}
\begin{itemize}
    \item Quilolitro $(Kl) = 1000 l = 10^3 l$
    \item Hectolitro $(hl) = 100 l = 10^3 l$
    \item Decalitro $(dal) = 10 l = 10^1l$
\end{itemize}
\subsection{Submúltiplos de Litro:}
\begin{itemize}
    \item Decilitro $(dl) = 0,1 l = 10^{-1}l$
    \item Centilitro $(cl) = 0,01 l = 10^{-2}l$
    \item Mililitro $(ml) = 0,001 l = 10^{-3}l$
\end{itemize}
\section{Medidas de Massa}
As medidas de massa são utilizadas para medir a quantidade
de massa em um corpo. No SI, a unidade fundamental de massa é
o quilograma (kg), embora o grama (g) seja a unidade mais
conveniente para a maioria das medidas de laboratório.
\subsection{Múltiplos de Grama:}
\begin{itemize}
    \item Quilograma $(Kg) = 1000 g = 10^3 g$
    \item Hectograma $(hg) = 100 g = 10^2 g$
    \item Decagrama $(dag) = 10 g = 10^1 g$
\end{itemize}
\subsection{Submúltiplos de Grama:}
\begin{itemize}
    \item Decigrama $(dg) = 0,1 g = 10^{-1} g$
    \item Centigrama $(cg) = 0,01 g = 10^{-2} g$
    \item Miligrama $(mg) = 0,001 g = 10^{-3}$
\end{itemize}
Existe ainda a unidade especial:
Tonelada $(t) = 1000 kg = 10^3 kg = 10^6 g$

\section{Exercícios}

Esta é um exemplo de exercício.
\begin{exercise}
Este é um bom lugar para fazer uma pergunta para testar o progresso da aprendizagem ou cimentar mais as idéias na mente dos alunos.
\end{exercise}


\chapter{Proporcionalidade}
Este tópico é fundamental, não apenas em Matemática, como na Matemática Financeira, na Física  na Química.Toda vez que trabalhamos com comparação de grandezas, algo comum na Física por exemplo, estamos utilizando os conceitos básicos da proporção.   

\section{Razão}
Razão pode ser definida como o quociente (divisão ou fração) entre dois números reais, sendo o quociente do primeiro número com o segundo número. Desta forma, a razão entre os números $"a"$ e $"b"$ será o quociente:

\begin{center}
	$\frac{a}{b}$ ou $a:b$, para $(b \neq 0)$. 	Sendo lida da seguinte forma: 
	$"a"$ está para $"b"$
\end{center}

\begin{example}
A razão $\frac{2}{3}$ é lida da seguinte forma: 2 está para 3.
\end{example}

\noindent Ou seja, basicamente vamos ter numa razão a divisão entre um número real $"a"$ e um número real $"b"$ que resultará em um número real $"c"$. Isto é:  

\begin{center}
	$\frac{a}{b} = c$ 
\end{center}

\begin{example}
   $\frac{3}{2} = 1,5$
\end{example}

\section{Proporção}

Proporção é definida como a igualdade entre duas ou mais razões. Isto é:

\begin{center}
    $\frac{a}{b} = \frac{c}{d}$, para $(b$ e $d \neq 0)$
\end{center}

Onde $"a"$ e $"d"$ são denominados extremos da proporção , $"b"$ e $"c"$ são os meios da proporção

\section{Escalas}

\begin{itemize}
	\item Escala
	\subitem Definição de escala 
	\subitem Esvala de apliação
	\subitem Escala de redução
\end{itemize}

\section{Porcentagem}
\begin{itemize}
	\item Razão centezimal
	\item Fator de multiplicação
	\item Interpretação situacional
\end{itemize}

\section{Exercícios}

Esta é um exemplo de exercício.

\begin{exercise}
Este é um bom lugar para fazer uma pergunta para testar o progresso da aprendizagem ou cimentar mais as idéias na mente dos alunos.
\end{exercise}


\part{Geometria}
\chapter{Geometria Plana}
    \section{Classificação dos polígonos}
    
    \section{Área dos polígonos}
    
    \section{Exercícios}
        Esta é um exemplo de exercício.
        \begin{exercise}
        Este é um bom lugar para fazer uma pergunta para testar o progresso da aprendizagem ou cimentar mais as idéias na mente dos alunos.
        \end{exercise} 
        
\chapter{Geometria Espacial}
    
\chapter{Geometria Analítica}


\part{Funções}
\chapter{Introdução}

\chapter{Função injetora ou injetiva}

\chapter{Função sobrejetora ou sobrejetiva}

\chapter{Função bijetora ou bijetiva}

\chapter{Função constante}

\chapter{Função par}

\chapter{Função ímpar}

\chapter{Função afim ou polinomial do primeiro grau}
    \begin{itemize}
    	\item O que é uma função
    	\item Domínio e imagem de uma função
    	\item Obtenção do domínio
    	\item Construção do gráfico cartesiano
    	\item Raízes de uma função
    	\item Propriedades de uma função
    	\item Função par e função ímpar
    	\item Funções crescente e decrescente
    	\item Função composta
    	\item Função inversa
    \end{itemize}

\chapter{Função linear}

\chapter{Função crescente}

\chapter{Função decrescente}

\chapter{Função quadrática ou polinomial do segundo grau}

\chapter{Função modular}

\chapter{Função exponencial}

\chapter{Função logarítmica}

\chapter{Função trigonométrica}

\chapter{Função raiz}


\section{Gráficos e Tabelas}% OK
\begin{itemize}
	\item Leitura e interpretação de gráficos (sistema ortogonal)
\end{itemize}

\section{Estatística}%ok
\begin{itemize}
	\item Rol
	\item Classe
	\item Moda, Média e mediana 
\end{itemize}

\section{Probabilidade}%ok
\begin{itemize}
	\item Espaço Amostral
	\item Eventos (seus tipos)
	\item Conceito de Probabilidade
	\item Probabilidade Condicional
	\item Multiplicação de Probabilidade
	\item Fórmulas.
\end{itemize}

\section{Trigonometria no Triângulo Retângulo}

\begin{itemize}
	\item   Definições de um triangulo retangulo
	\item razões trigonometricas
	\item Lei dos senos
\end{itemize}

\section{Análise Combinatória}%ok
\begin{itemize}
	\item Fatorial
	\item Arranjo simples
	\item Permutação simples
	\item Combinação simples
\end{itemize}

\section{Geometria Analítica}%ok
\begin{itemize}
	\item Fatorial
	\item Arranjo simples
	\item Permutação simples
	\item Combinação simples
\end{itemize}

\begin{center}
\textbf{EXERCÍCIOS}
    
\end{center}
	\begin{itemize}
		\item Fatorial
		\item Arranjo simples
		\item Permutação simples
		\item Combinação simples
	\end{itemize}
\end{document}
